% \section{Two Reinforcement Learning Paradigms for Option Pricing}
\section{Two RL Formulations for Option Pricing and Hedging}
\label{method}

% To overcome the limitations of the original QLBS formulation, two complementary strategies have been proposed. 
% The first approach extends the existing QLBS framework by redefining the value function $V_t^\pi$ as an $\{\cF_t\}$-adapted process.
% The second approach adopts a conceptually different perspective by constructing an adaptive portfolio value process and defining the reward in terms of the terminal hedging tracking error. 
% Both approaches explicitly account for transaction costs and are compatible with both value-based and policy-based reinforcement learning algorithms.

% To address the limitations of the original QLBS formulation, two complementary approaches can be pursued.
We pursue two complementary methods where
the first extends the QLBS framework by redefining the value function $V_t^\pi$ adapted to ${\cF_t}$,
and the second takes a forward viewpoint by constructing an adaptive portfolio value process and specifying the reward via the terminal hedging tracking error.
Both approaches incorporate transaction costs explicitly and remain compatible with both value-based and policy-based reinforcement learning algorithms.


\subsection{Adaptive-QLBS: Backward Value-Based RL}

\begin{defn}
\label{def:qlbs}
Let $d_T(t):=\left(1-\frac{t}{T}\right)$ and $\gamma:=e^{-r \Delta t}$.
The value function of the adaptive-QLBS method reads
{\small%
\begin{align}
% V_{t}^{\pi}\left(X_{t}\right)&:=\bE_{t}^{\pi}\left[-d_T(t)\Pi_{t}\left(X_{t}\right)-\lambda\sum_{\tau=t}^{T}\gamma^{\tau-t}\sqrt{\text{Var}\left[\Pi_{\tau}\left(X_{\tau}\right)\right]}\right],\label{eq:qlbs-modified-V}
V_{t}^{\pi}\left(X_{t}\right)&:=\bE_{t}^{\pi}\left[-d_T(t)\Pi_{t}-\lambda\sum_{\tau=t}^{T}\gamma^{\tau-t}\sqrt{\text{Var}\Pi_{\tau}}\right],\label{eq:qlbs-modified-V}
\end{align}
}with $R_{t+1}\left(X_{t},a_{t}\right):=V_t^\pi(X_t)-\bE_{t}^{\pi} V_{t+1}^\pi(X_{t+1})$ as the reward function where $X_{t+1}$ is implicitly determined by $X_t,a_t$ and the self-financing condition \cref{eq:self-financing}
\end{defn}

% Our modifications are twofold:
% we introduce a diminishing factor $d_T(t)$ that weights the portfolio term from $1$ at $t=0$ to $0$ at $t=T$, smoothing the contribution of the terminal payoff,
% and we replace variance terms by their square roots to obtain a dimensionless and numerically more stable value estimate.
% When transaction costs are included, the portfolio value process $\Pi_{t}$ is computed backward using the self-financing condition \cref{eq:self-financing}.
% A schematic illustration of the adaptive-QLBS model is provided in \cref{fig:qlbs-illustration}.

Our modifications are twofold:
the discounting factor $d_T(t)$ lessens the influence of the terminal payoff on the portfolio term,
and the variance terms are square rooted to obtain a dimensionless, numerically more stable value estimate.
With transaction costs, the portfolio value process $\Pi_t$ is computed backward via the self-financing relation in \cref{eq:self-financing}.
One can show that the estimated price $C(S_0)$ is monotonically increasing in both $\lambda$ and $\epsilon$ when $\epsilon\gg 1$.
% we introduce a discounting factor $d_T(t)$ that decreases from $1$ at $t=0$ to $0$ at $t=T$, thereby smoothing the influence of the terminal payoff on the portfolio term,
% and we replace variance terms with their square roots to obtain a dimensionless, numerically more stable value estimate.
% With transaction costs, the portfolio value process $\Pi_t$ is computed backward via the self-financing relation in \cref{eq:self-financing}.
A schematic overview of the adaptive-QLBS model is shown in \cref{fig:qlbs-illustration}.


\begin{figure*}[htbp]
    % \centering
    % \hfill
    \begin{minipage}{0.35\textwidth}
        \centering
        \begin{tikzpicture}[x=0.8cm,y=0.8cm, line cap=round, line join=round]
\begin{scope}[rotate=90]
  % ---- key coordinates ----
  \coordinate (O)   at (3.5,0); 
  \coordinate (t0)  at (4,0); \coordinate (t0a)  at (4,2.5);
  \coordinate (T)   at (7.5,0);
  \coordinate (tau) at (7.65,0);

  % ---- timeline (bottom) ----
  % past: solid
  \draw[thick,dotted] ($(O)+(0, -1)$) -- ($(t0)+(0, -1)$);
  % future: dotted
  \draw[thick,-stealth] ($(t0)+(0, -1)$) --($(T)+(0, -1)$);

  \node[right] at ($(O)+(0, -1)$) {$0$};
  \node[right]      at ($(t0)+(0, -1)$) {$t$};
  \node[right]      at ($(T)+(0, -1)$) {$T$};

  % ---- current time marker ----
  \fill ($(t0)+(0, -1)$) circle (2.2pt);

  % ---- state annotation (above timeline) ----
  % \node[above] at ($(t0a)+(0,1.1)$) {};


  % ---- scenario fan (all above timeline) ----
  \node[above] at ($(6.85,2.05)$) {$S_T$};
  \draw[color=gray] (t0a) --
  ++(0.055,0.117) -- ++(0.055,0.074) -- ++(0.055,0.149) -- ++(0.055,0.151) -- ++(0.055,0.152) -- ++(0.055,0.084) -- ++(0.055,0.018) -- ++(0.055,0.056) -- ++(0.055,0.105) -- ++(0.055,0.045) -- ++(0.055,0.083) -- ++(0.055,0.094) -- ++(0.055,0.085) -- ++(0.055,-0.047) -- ++(0.055,0.020) -- ++(0.055,0.103) -- ++(0.055,0.035) -- ++(0.055,0.006) -- ++(0.055,0.015) -- ++(0.055,0.013) -- ++(0.055,0.089) -- ++(0.055,-0.042) -- ++(0.055,0.073) -- ++(0.055,0.009) -- ++(0.055,0.006) -- ++(0.055,0.007) -- ++(0.055,0.054) -- ++(0.055,0.057) -- ++(0.055,0.050) -- ++(0.055,-0.035) -- ++(0.055,0.075) -- ++(0.055,0.051) -- ++(0.055,0.046) -- ++(0.055,-0.031) -- ++(0.055,0.088) -- ++(0.055,0.047) -- ++(0.055,0.019) -- ++(0.055,0.099) -- ++(0.055,-0.010) -- ++(0.055,0.073) -- ++(0.055,0.006) -- ++(0.055,-0.072) -- ++(0.055,0.027) -- ++(0.055,0.074) -- ++(0.055,-0.008) -- ++(0.055,-0.019) -- ++(0.055,-0.064) -- ++(0.055,0.020) -- ++(0.055,-0.035) -- ++(0.055,-0.007) -- ++(0.055,0.022) -- ++(0.055,-0.052) -- ++(0.055,0.058) -- ++(0.055,-0.052) -- ++(0.055,-0.032) -- ++(0.055,0.067) -- ++(0.055,-0.109) -- ++(0.055,0.025) -- ++(0.055,-0.021) -- ++(0.055,0.065) -- ++(0.055,0.026) -- ++(0.055,0.067) -- ++(0.055,-0.016) ;
  \draw[color=gray] (t0a) --
  ++(0.055,0.049) -- ++(0.055,0.082) -- ++(0.055,0.053) -- ++(0.055,0.020) -- ++(0.055,0.061) -- ++(0.055,0.102) -- ++(0.055,0.079) -- ++(0.055,-0.002) -- ++(0.055,-0.029) -- ++(0.055,0.003) -- ++(0.055,0.034) -- ++(0.055,-0.079) -- ++(0.055,0.024) -- ++(0.055,-0.025) -- ++(0.055,0.000) -- ++(0.055,0.009) -- ++(0.055,0.019) -- ++(0.055,0.052) -- ++(0.055,0.079) -- ++(0.055,0.020) -- ++(0.055,0.089) -- ++(0.055,-0.011) -- ++(0.055,0.037) -- ++(0.055,0.061) -- ++(0.055,0.019) -- ++(0.055,-0.021) -- ++(0.055,-0.028) -- ++(0.055,-0.005) -- ++(0.055,0.027) -- ++(0.055,-0.032) -- ++(0.055,0.007) -- ++(0.055,0.009) -- ++(0.055,0.042) -- ++(0.055,0.024) -- ++(0.055,0.029) -- ++(0.055,-0.020) -- ++(0.055,0.006) -- ++(0.055,0.049) -- ++(0.055,0.079) -- ++(0.055,-0.054) -- ++(0.055,0.079) -- ++(0.055,0.067) -- ++(0.055,0.036) -- ++(0.055,0.010) -- ++(0.055,-0.018) -- ++(0.055,0.067) -- ++(0.055,0.086) -- ++(0.055,0.074) -- ++(0.055,0.047) -- ++(0.055,-0.000) -- ++(0.055,-0.074) -- ++(0.055,-0.013) -- ++(0.055,0.019) -- ++(0.055,-0.074) -- ++(0.055,0.009) -- ++(0.055,0.011) -- ++(0.055,0.022) -- ++(0.055,-0.067) -- ++(0.055,-0.039) -- ++(0.055,-0.026) -- ++(0.055,0.026) -- ++(0.055,0.067) -- ++(0.055,-0.016) ;
  \draw[thick] (t0a) --
  ++(0.055,-0.009) -- ++(0.055,0.007) -- ++(0.055,-0.004) -- ++(0.055,-0.056) -- ++(0.055,0.029) -- ++(0.055,-0.023) -- ++(0.055,0.013) -- ++(0.055,-0.030) -- ++(0.055,0.054) -- ++(0.055,-0.011) -- ++(0.055,0.011) -- ++(0.055,-0.004) -- ++(0.055,0.023) -- ++(0.055,-0.002) -- ++(0.055,0.004) -- ++(0.055,0.038) -- ++(0.055,0.009) -- ++(0.055,0.021) -- ++(0.055,-0.024) -- ++(0.055,0.021) -- ++(0.055,-0.030) -- ++(0.055,0.020) -- ++(0.055,-0.045) -- ++(0.055,-0.050) -- ++(0.055,-0.009) -- ++(0.055,0.045) -- ++(0.055,0.008) -- ++(0.055,-0.015) -- ++(0.055,-0.004) -- ++(0.055,-0.065) -- ++(0.055,0.017) -- ++(0.055,-0.029) -- ++(0.055,0.022) -- ++(0.055,0.057) -- ++(0.055,0.005) -- ++(0.055,-0.007) -- ++(0.055,0.029) -- ++(0.055,0.010) -- ++(0.055,-0.012) -- ++(0.055,-0.031) -- ++(0.055,-0.043) -- ++(0.055,0.048) -- ++(0.055,0.019) -- ++(0.055,0.061) -- ++(0.055,-0.011) -- ++(0.055,-0.032) -- ++(0.055,0.030) -- ++(0.055,-0.026) -- ++(0.055,-0.058) -- ++(0.055,0.049) -- ++(0.055,-0.026) -- ++(0.055,0.014) -- ++(0.055,-0.019) -- ++(0.055,0.037) -- ++(0.055,-0.025) -- ++(0.055,0.009) -- ++(0.055,0.005) -- ++(0.055,-0.005) -- ++(0.055,0.044) -- ++(0.055,0.026) -- ++(0.055,0.026) -- ++(0.055,0.067);
  \draw[color=gray] (t0a) --
  ++(0.055,-0.127) -- ++(0.055,0.040) -- ++(0.055,-0.173) -- ++(0.055,-0.052) -- ++(0.055,-0.075) -- ++(0.055,-0.066) -- ++(0.055,-0.064) -- ++(0.055,-0.034) -- ++(0.055,0.024) -- ++(0.055,0.034) -- ++(0.055,-0.066) -- ++(0.055,-0.038) -- ++(0.055,-0.038) -- ++(0.055,-0.013) -- ++(0.055,-0.014) -- ++(0.055,0.023) -- ++(0.055,-0.013) -- ++(0.055,0.066) -- ++(0.055,-0.027) -- ++(0.055,-0.012) -- ++(0.055,-0.045) -- ++(0.055,-0.074) -- ++(0.055,-0.084) -- ++(0.055,0.018) -- ++(0.055,-0.079) -- ++(0.055,-0.042) -- ++(0.055,-0.057) -- ++(0.055,-0.009) -- ++(0.055,0.005) -- ++(0.055,-0.062) -- ++(0.055,-0.036) -- ++(0.055,0.014) -- ++(0.055,-0.038) -- ++(0.055,-0.035) -- ++(0.055,-0.006) -- ++(0.055,-0.051) -- ++(0.055,0.176) -- ++(0.055,0.042) -- ++(0.055,-0.060) -- ++(0.055,0.025) -- ++(0.055,0.002) -- ++(0.055,-0.044) -- ++(0.055,0.036) -- ++(0.055,0.027) -- ++(0.055,-0.097) -- ++(0.055,0.029) -- ++(0.055,-0.003) -- ++(0.055,-0.066) -- ++(0.055,-0.080) -- ++(0.055,-0.048) -- ++(0.055,-0.035) -- ++(0.055,0.035) -- ++(0.055,0.086) -- ++(0.055,0.015) -- ++(0.055,0.064) -- ++(0.055,0.049) -- ++(0.055,0.046) -- ++(0.055,-0.136) -- ++(0.055,0.072) -- ++(0.055,0.016)  -- ++(0.055,0.026) -- ++(0.055,0.067);
  \draw[color=gray] (t0a) --
  ++(0.055,-0.032) -- ++(0.055,-0.053) -- ++(0.055,-0.214) -- ++(0.055,-0.108) -- ++(0.055,-0.126) -- ++(0.055,-0.151) -- ++(0.055,-0.081) -- ++(0.055,-0.168) -- ++(0.055,-0.102) -- ++(0.055,-0.070) -- ++(0.055,-0.016) -- ++(0.055,-0.028) -- ++(0.055,-0.087) -- ++(0.055,-0.023) -- ++(0.055,-0.027) -- ++(0.055,-0.027) -- ++(0.055,-0.137) -- ++(0.055,0.060) -- ++(0.055,-0.097) -- ++(0.055,-0.146) -- ++(0.055,-0.040) -- ++(0.055,-0.066) -- ++(0.055,-0.011) -- ++(0.055,0.046) -- ++(0.055,-0.061) -- ++(0.055,-0.062) -- ++(0.055,0.002) -- ++(0.055,0.001) -- ++(0.055,0.013) -- ++(0.055,0.071) -- ++(0.055,0.059) -- ++(0.055,0.001) -- ++(0.055,-0.013) -- ++(0.055,0.016) -- ++(0.055,0.022) -- ++(0.055,0.031) -- ++(0.055,0.022) -- ++(0.055,-0.054) -- ++(0.055,0.021) -- ++(0.055,-0.017) -- ++(0.055,-0.012) -- ++(0.055,-0.009) -- ++(0.055,-0.011) -- ++(0.055,0.031) -- ++(0.055,0.056) -- ++(0.055,-0.114) -- ++(0.055,-0.057) -- ++(0.055,0.092) -- ++(0.055,-0.048) -- ++(0.055,0.058) -- ++(0.055,-0.019) -- ++(0.055,0.110) -- ++(0.055,0.034) -- ++(0.055,-0.016) -- ++(0.055,-0.015) -- ++(0.055,-0.041) -- ++(0.055,-0.065) -- ++(0.055,-0.102) -- ++(0.055,-0.056) -- ++(0.055,-0.058) -- ++(0.055,-0.028) -- ++(0.055,-0.034) ;
  \draw[thick,dotted] (t0a) --
  ++(-0.050,0.049) -- ++(-0.050,-0.165) -- ++(-0.050,0.030) -- ++(-0.050,0.037) -- ++(-0.050,0.257) -- ++(-0.050,-0.076) -- ++(-0.050,-0.015) -- ++(-0.050,0.020) -- ++(-0.050,-0.223) -- ++(-0.050,-0.071);
  %-- ++(-0.050,0.079) -- ++(-0.050,-0.094) -- ++(-0.050,-0.127) -- ++(-0.050,0.003) -- ++(-0.050,0.103) -- ++(-0.050,-0.191) -- ++(-0.050,-0.141) -- ++(-0.050,-0.083) -- ++(-0.050,-0.069) -- ++(-0.050,0.002) -- ++(-0.050,-0.159) -- ++(-0.050,0.109) -- ++(-0.050,0.158) -- ++(-0.050,0.006) -- ++(-0.050,-0.092) -- ++(-0.050,0.014);

  \draw[-{Stealth[length=2mm,width=1.5mm]}] ($(t0)+(0,1.5)$) -- ($(t0)+(0,2)$);

  % ---- cross-section for dispersion ----
  \draw[dashed,thick] ($(tau)+(0,0.3)$) -- ($(tau)+(0,4.8)$);

  % \draw[decorate, decoration={brace, amplitude=6pt}, thick]
  %   ($(tau)+(0,4.6)$) -- ($(tau)+(0,0.6)$);

  % labels for dispersion and variance term
  \node[align=center] at ($(t0a)+(1.3,2.7)$)
    {$\text{Var}\,\Pi_t$ \\under\\ \cref{eq:self-financing}};

    
  \node[align=center,above] at ($(t0a)+(-0.7,-2.05)$) {hedge $u_t$ \\ state $(t,S_t)$};

  \node[align=right] at ($(tau)+(0.5,2.5)$)
    {terminal condition $h$};

\end{scope}
\end{tikzpicture}

        % \includegraphics[width=\linewidth]{backup/qlbs-illustration}
        
        \label{fig:qlbs-illustration}
    \end{minipage}
    % \hfill
    \hspace{0.02\textwidth}
    % Second figure
    \begin{minipage}{0.62\textwidth}
        \centering
        \begin{tikzpicture}[x=0.8cm,y=0.8cm,>=Stealth]
\begin{scope}[yscale=-1]
  % ---------- styles ----------
  \tikzset{
    tl/.style={line width=0.9pt},
    pf/.style={line width=0.9pt},
    tick/.style={line width=0.9pt},
    darr/.style={dashed, line width=0.8pt, ->},
    dot/.style={circle, fill=black, inner sep=1.3pt},
  }

  % ---------- key x-positions ----------
  % Use symbolic positions for readability; adjust if you want different spacing
  \def\xO{0}
  \def\xOne{1}
  \def\xTwo{2}
  \def\xTmid{6}  % position of t on the top timeline
  \def\xT{9}     % position of T on the top timeline

  % ---------- top timeline ----------
  \draw[tl] (\xO,0) -- (\xT,0);

  % ticks + labels on timeline
  \foreach \x/\lab in {\xO/0,\xOne/1,\xTwo/2,\xTmid/$t$,\xT/$T$}{
    \draw[tick] (\x,0.12) -- (\x,-0.12);
    \node[below] at (\x,0.18) {\lab};
  }
  \node[below] at (3.9,0.2) {$\cdots$};
  \node[below] at (7.5,0.2) {$\cdots$};

  % ---------- left labels ----------
  \node[anchor=east] at (-0.25,-0.3) {portfolio};
  \node[anchor=east,color=gray] at (-0.25,-0.9) {\#1};
  % \node[anchor=east,color=gray] at (-0.25,-2.00) {\#2};
  \node[anchor=east] at (-0.25,-2.00) {\#$t$};
  \node[anchor=east,color=gray] at (-0.25,-3.3) {\#$T$};

  % ---------- portfolio #1 (0 -> 1) ----------
  \node[dot,color=gray] at (\xO,-0.90) {};
  \draw[pf,color=gray] (\xO,-0.90) -- (\xOne,-0.90);
  % bracket indicating u_0^(1)
  % \draw[line width=0.9pt] (\xO+0.05,-1.52) -- (\xOne-0.05,-1.52);
  % \draw[line width=0.9pt] (\xO+0.05,-1.52) -- (\xO+0.05,-1.42);
  % \draw[line width=0.9pt] (\xOne-0.05,-1.52) -- (\xOne-0.05,-1.42);
  \node[above,color=gray] at (\xO+0.5,-0.90) {$u^{(1)}_{0}$};

  % reward arrow R1 at time 1
  \draw[darr,color=gray] (\xOne,-0.90) -- (\xOne,-0.2);
  \node[right,color=gray] at (\xOne, -0.7) {$R_{1}$};

  % ---------- portfolio #2 (0 -> 2) ----------
  % \node[dot,color=gray] at (\xO,-2.00) {};
  % \draw[pf,color=gray] (\xO,-2.00) -- (\xTwo,-2.00);
  % \draw[tick,color=gray] (\xOne,-2.00+0.12) -- (\xOne,-2.00-0.12);
  % \node[above,color=gray] at (\xO+0.5,-2.00) {$u^{(2)}_{0}$};
  % \node[above,color=gray] at (\xOne+0.50,-2.00) {$u^{(2)}_{1}$};

  % % reward arrow R2 at time 2
  % \draw[darr,color=gray] (\xTwo,-2.00) -- (\xTwo,-0.2);
  % \node[right,color=gray] at (\xTwo, -1.10) {$R_{2}$};

  % % vertical ellipsis between #2 and #t
  \node[color=gray] at (-0.55,-1.55) {$\vdots$};

  % ---------- portfolio #t (0 -> t) ----------
  \node[dot] at (\xO,-2.0) {};
  \draw[pf] (\xO,-2.0) -- (\xTmid,-2.0);
  \draw[tick] (\xOne,-2.0+0.12) -- (\xOne,-2.0-0.12);
  \draw[tick] (\xOne+1.0,-2.0+0.12) -- (\xOne+1.0,-2.0-0.12);
  \draw[tick] (\xTmid-1.0,-2.0+0.12) -- (\xTmid-1.0,-2.0-0.12);
  \node[above] at (\xO+0.5,-2.0) {$u^{(t)}_{0}$};
  \node[above] at (\xOne+0.50,-2.0) {$u^{(t)}_{1}$};
  \node[above] at (\xTmid-0.5,-2.0) {$u^{(t)}_{t-1}$};
  \node[above] at (3.9,-2.0) {$\cdots$};

  % reward arrow Rt at time t
  \draw[darr] (\xTmid,-2.0) -- (\xTmid,-0.2);
  \node[right] at (\xTmid, -0.9) {$R_{t}$};

  % vertical ellipsis between #t and #T
  \node[color=gray] at (-0.55,-2.70) {$\vdots$};

  % ---------- portfolio #T (0 -> T) ----------
  \node[dot,color=gray] at (\xO,-3.3) {};
  \draw[pf,color=gray] (\xO,-3.3) -- (\xT,-3.3);
  \draw[tick,color=gray] (\xOne,-3.3+0.12) -- (\xOne,-3.3-0.12);
  \draw[tick,color=gray] (\xOne+1,-3.3+0.12) -- (\xOne+1,-3.3-0.12);
  \draw[tick,color=gray] (\xT-1,-3.3+0.12) -- (\xT-1,-3.3-0.12);
  \node[above,color=gray] at (\xO+0.5,-3.3) {$u^{(T)}_{0}$};
  \node[above,color=gray] at (\xOne+0.50,-3.3) {$u^{(T)}_{1}$};
  \node[above,color=gray] at (\xT-0.5,-3.3) {$u^{(T)}_{T-1}$};
  \node[above,color=gray] at (5.0,-3.3) {$\cdots$};

  % reward arrow RT at time T
  \draw[darr,color=gray] (\xT,-3.3) -- (\xT,-0.2);
  \node[right,color=gray] at (\xT, -1.5) {$R_{T}$};
\end{scope}
\end{tikzpicture}

        % \includegraphics[width=\linewidth]{backup/rlop-illustration.png}
        % \label{fig:rlop-illustration}
    \end{minipage}
    % \hfill
    
    \caption{Adaptive-QLBS (left, backward \& value-based) vs RLOP (right, forward \& replication-based).}

\end{figure*}


% The proposition below explains why the option price is increasing in the key parameters where higher risk aversion 
% $\lambda$ and larger transaction friction both raise the price. We first prove monotonicity of the value function under any fixed policy, and then obtain monotonicity of the optimal value by maximization. 



% The proposition below clarifies why the option price increases with the key parameters: greater risk aversion $\lambda$ and higher transaction frictions both lead to a higher price. We first establish monotonicity of the value function under any fixed policy, and then derive monotonicity of the optimal value by taking the maximizer.

% \begin{prop}
For sufficiently large $\epsilon$ that appears in the linear transaction cost assumption $\text{TC}(\Delta u, S)=\epsilon|\Delta u| \,S$, the option price $C(S_0) := -\max_{\pi\in\mathbf{\Pi}} V_0^\pi$ is monotonically increasing in both $\lambda$ and $\epsilon$.
\label{prop:qlbs-monotone}
\end{prop}
\begin{proof}
% Given any fixed policy $\pi$ and under the self-financing condition \cref{eq:self-financing}, and the assumed linear transaction cost (??), the portfolio value at time $t$ may be written as
Given any policy $\pi$, the portfolio value at time $t$ may be written as
\begin{align}
\Pi_t &= u_t  S_t + \epsilon \mathfrak{S}_t + \gamma^{T-t} \mathfrak{E}_t, \label{eq:qlbs-proof-port-value} \\
\mathfrak{S}_t :&= \sum_{j=0}^{T-t-1} \gamma^{j+1} |\Delta u_{t+j}| S_{t+j+1}, \nonumber \\
\mathfrak{E}_t :&= h(S_T) - u_{T-1}S_T \nonumber
\end{align}
under the self-financing condition
% \cref{eq:self-financing}
(eq. 1)
.
This expression is affine in $\epsilon$ and independent of $\lambda$.
Hence, the monotonicity of $V_t^\pi$ with respect to  $\lambda$ follows immediately from
% \cref{eq:qlbs-modified-V}
(eq. 2)
, since the discounted risk term is always non-positive.
For $\epsilon\gg1$, the monotonicity in $\epsilon$ follows from the fact that $\text{Var}_t\Pi_{t}\left(X_{t}\right)$ is a quadratic function of $\epsilon$ with non-negative leading coefficient.
% which is independent of $\lambda$ and affine in $\epsilon$. Thus, the monitonicity of $V_t^\pi$ with respect to $\lambda$ immediately follows since the discounted risk in \ref{eq:qlbs-modified-V} is always non-positive.
% The monotonicty to $\epsilon$ in large $\epsilon$ regime is a natural consequence of $\text{Var}_t\Pi_{t}\left(X_{t}\right)$ being a quadratic form in $\epsilon$.

To establish that $C(S_0,\lambda)$ is increasing in $\lambda$, let $\pi(\lambda)$ denote the maximizer of $V_t^\pi$.
For $\lambda ' > \lambda$, we have
$V_t^{\pi(\lambda )}(S_t;\lambda) \ge V_t^{\pi(\lambda')}(S_t;\lambda)$
because $\pi(\lambda)$  maximizes the value function at level $\lambda$.
From the argument above, $V_t^{\pi(\lambda')}(S_t;\lambda') \le V_t^{\pi(\lambda')}(S_t;\lambda)$.
Combining these inequalities yields
\begin{align*}
C(S_0,\lambda') &= -V_0^{\pi(\lambda')} (S_0;\lambda') \ge -V_0^{\pi(\lambda')} (S_0;\lambda) \\
&\ge -V_0^{\pi(\lambda)} (S_0;\lambda) = C(S_0,\lambda) .
\end{align*}
Thus, $C$ is monotone in $\lambda$.
The argument for its monotonicity in $\epsilon$ proceeds similarly.

% To show that $C(S_0,\lambda)$ is increasing in $\lambda$, let us introduce the maximizer $\pi(\lambda)$ for $V_t^\pi$. Then, given $\lambda ' > \lambda$, $V_t^{\pi(\lambda )}(S_t;\lambda) \ge V_t^{\pi(\lambda')}(S_t;\lambda)$ due to $\pi(\lambda )$ being $V_t^{\cdot}(S_t;\lambda)$'s maximier. We also establish that 
% $
% V_t^{\pi(\lambda')}(S_t;\lambda') \le V_t^{\pi(\lambda')}(S_t;\lambda)
% $
% in the paragraph above.
% \begin{align}
% C(S_0,\lambda') = -V_0^{\pi(\lambda')} (S_0;\lambda') \ge -V_0^{\pi(\lambda')} (S_0;\lambda) \ge -V_0^{\pi(\lambda)} (S_0;\lambda) = C(S_0,\lambda) .
% \end{align}
% The monitonicty of $C$ to $\epsilon$ can be established in a similar fashion.

% As for the monotonicty to $\epsilon$, notice that $\partial_\epsilon \Pi_t = \mathfrak{S}_t \ge 0$, it suffices to shown that the risk component $\text{Var}_t\Pi_{t}\left(X_{t}\right)$ is increasingly in $\epsilon$.
% To see this, we expand the variance by
% \begin{equation}
% \text{Var}_t\Pi_{t}\left(X_{t}\right) = \text{Var}_t\left( \gamma^{T-t} \mathfrak{E}_t\right) + \epsilon^2 \text{Var}_t\left(\mathfrak{S}_t\right) + \epsilon \text{Cov}_t\left(\gamma^{T-t} \mathfrak{E}_t, \mathfrak{S}_t\right) 
% \label{eq:qlbd-proof-expand-var}
% \end{equation}
\end{proof}
